% GENERATED FILE. Edit canonical lessons, then run npm run generate:books.
% canonical-source-hash: fnv1a64:dffdd47cf15aefc4
% canonical-lessons: TA-C459-anantam, TA-C459-inpam, TA-C459-purippu, TA-C459-kutukalam, TA-C459-ullacam

\chapter{Joy and High Spirits}
\label{ch:ta-joy-and-high-spirits}
\hlchaptermodality{459}

\begin{chapteropening}
\textbf{By the end of this chapter:} \emph{I can talk about bliss, delight and high spirits.}

Complete the last lesson of chapter 459: I can talk about bliss, delight and high spirits.
\end{chapteropening}

\section[\texorpdfstring{\ta{ஆனந்தம்} (āṉantam)}{ஆனந்தம் (āṉantam)}]{\ta{ஆனந்தம்} (āṉantam) — bliss; great joy}

\label{lesson:TA-C459-anantam}



\begin{quote}
Before the new one: say the Tamil for to lament, then the Tamil for to scream.
\end{quote}

\subsection*{You'll want to know: \ta{ஆனந்தம்}}
\textbf{\ta{ஆனந்தம்}} — \emph{āṉantam} — \textquotedblleft{}bliss; great joy\textquotedblright{}.

\begin{quote}
\ta{பேரனைப்} \ta{பார்த்ததும்} \ta{தாத்தாவுக்கு} \ta{ஆனந்தம்}.
\end{quote}

\begin{quote}
\emph{pēraṉaip pārttatum tāttāvukku āṉantam.}
\end{quote}

\begin{quote}
Grandfather was overjoyed to see his grandson.
\end{quote}

\ta{ஆனந்தக்} \ta{கண்ணீர்}, tears of joy.

\begin{cousinweb}
\textbf{\ta{ஆனந்தம்}} (\emph{āṉantam}) is borrowed from Sanskrit \emph{ānanda}, \textquotedblleft{}happiness\textquotedblright{} (CDIAL 1172), with the Tamil ending -\ta{அம்}.
\end{cousinweb}

\begin{grammarlens}[title={What you've built}]
One more word for describing what you lived and felt.
\end{grammarlens}

\subsection*{Your turn}
\begin{itemize}
\raggedright
  \item \emph{Say it:} \emph{āṉantam}
  \item \emph{Say it:} \emph{āṉantam}, once more
  \item \emph{Say it:} \emph{pēraṉaip pārttatum tāttāvukku āṉantam.} — Grandfather was overjoyed to see his grandson.
  \item \emph{Recall:} say the Tamil for to lament, then the Tamil for to scream, then say \emph{āṉantam} again
\end{itemize}

\begin{tcolorbox}[breakable,colback=teal!4,colframe=teal!35!black,title={Before you move on}]
What does \ta{ஆனந்தம்} mean? (\textquotedblleft{}Bliss; great joy\textquotedblright{}.) Say it once more.
\end{tcolorbox}
\section[\texorpdfstring{\ta{இன்பம்} (iṉpam)}{இன்பம் (iṉpam)}]{\ta{இன்பம்} (iṉpam) — pleasure; delight}

\label{lesson:TA-C459-inpam}



\begin{quote}
Before the new one: say the Tamil for to scream, then the Tamil for bliss.
\end{quote}

\subsection*{You'll want to know: \ta{இன்பம்}}
\textbf{\ta{இன்பம்}} — \emph{iṉpam} — \textquotedblleft{}pleasure; delight\textquotedblright{}.

\begin{quote}
\ta{தோட்ட} \ta{வேலை} \ta{எனக்கு} \ta{இன்பம்} \ta{தருகிறது}.
\end{quote}

\begin{quote}
\emph{tōṭṭa vēlai eṉakku iṉpam tarukiṟatu.}
\end{quote}

\begin{quote}
Garden work gives me pleasure.
\end{quote}

\ta{இன்} \textquotedblleft{}sweet\textquotedblright{} + -\ta{பம்}; \ta{இனிமையான} \textquotedblleft{}pleasant\textquotedblright{} comes from the same root.

\begin{cousinweb}
\textbf{\ta{இன்பம்}} (\emph{iṉpam}) is native Dravidian, built on \ta{இன்} \emph{iṉ}, \textquotedblleft{}sweet, pleasant\textquotedblright{} (DEDR 530).
\end{cousinweb}

\begin{grammarlens}[title={What you've built}]
One more word for describing what you lived and felt.
\end{grammarlens}

\subsection*{Your turn}
\begin{itemize}
\raggedright
  \item \emph{Say it:} \emph{iṉpam}
  \item \emph{Say it:} \emph{iṉpam}, once more
  \item \emph{Say it:} \emph{tōṭṭa vēlai eṉakku iṉpam tarukiṟatu.} — Garden work gives me pleasure.
  \item \emph{Recall:} say the Tamil for to scream, then the Tamil for bliss, then say \emph{iṉpam} again
\end{itemize}

\begin{tcolorbox}[breakable,colback=teal!4,colframe=teal!35!black,title={Before you move on}]
What does \ta{இன்பம்} mean? (\textquotedblleft{}Pleasure; delight\textquotedblright{}.) Say it once more.
\end{tcolorbox}
\section[\texorpdfstring{\ta{பூரிப்பு} (pūrippu)}{பூரிப்பு (pūrippu)}]{\ta{பூரிப்பு} (pūrippu) — elation; swelling joy}

\label{lesson:TA-C459-purippu}



\begin{quote}
Before the new one: say the Tamil for bliss, then the Tamil for pleasure.
\end{quote}

\subsection*{You'll want to know: \ta{பூரிப்பு}}
\textbf{\ta{பூரிப்பு}} — \emph{pūrippu} — \textquotedblleft{}elation; swelling joy\textquotedblright{}.

\begin{quote}
\ta{மகள்} \ta{வென்றதைக்} \ta{கேட்டு} \ta{அம்மா} \ta{பூரிப்பு} \ta{அடைந்தாள்}.
\end{quote}

\begin{quote}
\emph{makaḷ veṉṟataik kēṭṭu ammā pūrippu aṭaintāḷ.}
\end{quote}

\begin{quote}
Hearing that her daughter had won, the mother swelled with joy.
\end{quote}

\ta{பூரிப்பு} \ta{அடை}, to be filled with joy.

\begin{cousinweb}
The origin of \textbf{\ta{பூரிப்பு}} (\emph{pūrippu}) is not settled; it is sometimes linked with Sanskrit words for filling, but our sources do not confirm the path.
\end{cousinweb}

\begin{grammarlens}[title={What you've built}]
One more word for describing what you lived and felt.
\end{grammarlens}

\subsection*{Your turn}
\begin{itemize}
\raggedright
  \item \emph{Say it:} \emph{pūrippu}
  \item \emph{Say it:} \emph{pūrippu}, once more
  \item \emph{Say it:} \emph{makaḷ veṉṟataik kēṭṭu ammā pūrippu aṭaintāḷ.} — Hearing that her daughter had won, the mother swelled with joy.
  \item \emph{Recall:} say the Tamil for bliss, then the Tamil for pleasure, then say \emph{pūrippu} again
\end{itemize}

\begin{tcolorbox}[breakable,colback=teal!4,colframe=teal!35!black,title={Before you move on}]
What does \ta{பூரிப்பு} mean? (\textquotedblleft{}Elation; swelling joy\textquotedblright{}.) Say it once more.
\end{tcolorbox}
\section[\texorpdfstring{\ta{குதூகலம்} (kutūkalam)}{குதூகலம் (kutūkalam)}]{\ta{குதூகலம்} (kutūkalam) — glee; high spirits}

\label{lesson:TA-C459-kutukalam}



\begin{quote}
Before the new one: say the Tamil for pleasure, then the Tamil for elation.
\end{quote}

\subsection*{You'll want to know: \ta{குதூகலம்}}
\textbf{\ta{குதூகலம்}} — \emph{kutūkalam} — \textquotedblleft{}glee; high spirits\textquotedblright{}.

\begin{quote}
\ta{விடுமுறை} \ta{வந்ததால்} \ta{குழந்தைகள்} \ta{குதூகலமாக} \ta{இருந்தார்கள்}.
\end{quote}

\begin{quote}
\emph{viṭumuṟai vantatāl kuḻantaikaḷ kutūkalamāka iruntārkaḷ.}
\end{quote}

\begin{quote}
The holidays had come, so the children were in high spirits.
\end{quote}

With -\ta{ஆக} and \ta{இரு}: \ta{குதூகலமாக} \ta{இருந்தார்கள்}, they were gleeful.

\begin{cousinweb}
\textbf{\ta{குதூகலம்}} (\emph{kutūkalam}) is borrowed from Sanskrit \emph{kutūhala}, \textquotedblleft{}curiosity, eager interest\textquotedblright{} (Monier-Williams); in Tamil it is high-spirited joy.
\end{cousinweb}

\begin{grammarlens}[title={What you've built}]
One more word for describing what you lived and felt.
\end{grammarlens}

\subsection*{Your turn}
\begin{itemize}
\raggedright
  \item \emph{Say it:} \emph{kutūkalam}
  \item \emph{Say it:} \emph{kutūkalam}, once more
  \item \emph{Say it:} \emph{viṭumuṟai vantatāl kuḻantaikaḷ kutūkalamāka iruntārkaḷ.} — The holidays had come, so the children were in high spirits.
  \item \emph{Recall:} say the Tamil for pleasure, then the Tamil for elation, then say \emph{kutūkalam} again
\end{itemize}

\begin{tcolorbox}[breakable,colback=teal!4,colframe=teal!35!black,title={Before you move on}]
What does \ta{குதூகலம்} mean? (\textquotedblleft{}Glee; high spirits\textquotedblright{}.) Say it once more.
\end{tcolorbox}
\section[\texorpdfstring{\ta{உல்லாசம்} (ullācam)}{உல்லாசம் (ullācam)}]{\ta{உல்லாசம்} (ullācam) — merriment; carefree fun}

\label{lesson:TA-C459-ullacam}



\begin{quote}
Before the new one: say the Tamil for elation, then the Tamil for glee.
\end{quote}

\subsection*{You'll want to know: \ta{உல்லாசம்}}
\textbf{\ta{உல்லாசம்}} — \emph{ullācam} — \textquotedblleft{}merriment; carefree fun\textquotedblright{}.

\begin{quote}
\ta{கடற்கரையில்} \ta{நாங்கள்} \ta{உல்லாசமாக} \ta{நேரம்} \ta{கழித்தோம்}.
\end{quote}

\begin{quote}
\emph{kaṭaṟkaraiyil nāṅkaḷ ullācamāka nēram kaḻittōm.}
\end{quote}

\begin{quote}
We spent a carefree time at the beach.
\end{quote}

\ta{உல்லாசப்} \ta{பயணம்}, a pleasure trip.

\begin{cousinweb}
\textbf{\ta{உல்லாசம்}} (\emph{ullācam}) is borrowed from Sanskrit \emph{ullāsa}, \textquotedblleft{}brightness; coming forth\textquotedblright{} (Monier-Williams); Tamil uses it for carefree enjoyment.
\end{cousinweb}

\begin{grammarlens}[title={What you've built}]
One more word for describing what you lived and felt.
\end{grammarlens}

\subsection*{Your turn}
\begin{itemize}
\raggedright
  \item \emph{Say it:} \emph{ullācam}
  \item \emph{Say it:} \emph{ullācam}, once more
  \item \emph{Say it:} \emph{kaṭaṟkaraiyil nāṅkaḷ ullācamāka nēram kaḻittōm.} — We spent a carefree time at the beach.
  \item \emph{Recall:} say the Tamil for elation, then the Tamil for glee, then say \emph{ullācam} again
\end{itemize}

\begin{tcolorbox}[breakable,colback=teal!4,colframe=teal!35!black,title={Before you move on}]
What does \ta{உல்லாசம்} mean? (\textquotedblleft{}Merriment; carefree fun\textquotedblright{}.) Say it once more.
\end{tcolorbox}
